% ============================================================
%  摘要
% ============================================================
\chapter*{摘\quad 要}
\addcontentsline{toc}{chapter}{摘要}

随着高校社团数量和活动规模持续增长，传统依赖纸质表单、线下审批与人工通知的管理方式逐渐暴露出流程分散、效率偏低和数据难以追溯等问题。针对这一场景，本文设计并实现了一套基于 Spring Boot 的校园社团管理系统。系统采用模块化单体架构，以 Spring Boot 3.x、Vue 3、MySQL 8、Redis 与 RabbitMQ 为核心技术栈，构建了用户认证、社团生命周期管理、招新、活动、公告通知、资源、财务以及 AI 问答与推荐等功能模块。在实现上，系统结合 JWT 与 HttpOnly Cookie 完成认证控制，利用 WebSocket 和异步通知提升关键消息触达效率，并引入基于 RAG 的知识库问答能力增强信息查询。测试结果表明，系统在当前单节点环境下已具备基本可运行性：关键回归用例复跑通过，登录、社团列表、活动报名与签到等核心链路在既定并发脚本下保持较低响应时间且错误率为 0\%。总体而言，该系统能够较好支撑高校社团管理的主要业务流程，但在长期稳定性、AI 评估完备性与更大规模部署验证方面仍需继续完善。

\vspace{1em}
\noindent\textbf{关键词：}Spring Boot；社团管理系统；JWT 认证；实时通知；RAG 问答

\newpage

% ============================================================
%  ABSTRACT
% ============================================================
\chapter*{ABSTRACT}
\addcontentsline{toc}{chapter}{ABSTRACT}

With the growth of student clubs and campus activities, traditional club management methods based on paper forms, offline approvals, and manual notifications have become increasingly inefficient and difficult to trace. To address this problem, this thesis designs and implements a campus club management system based on Spring Boot. The system adopts a modular monolithic architecture and uses Spring Boot 3.x, Vue 3, MySQL 8, Redis, and RabbitMQ as its core technology stack. It provides user authentication, club lifecycle management, recruitment, activity management, announcements and notifications, resource management, finance management, as well as AI-assisted question answering and recommendation functions.

In implementation, the system combines JWT and HttpOnly cookies for authentication control, uses WebSocket and asynchronous notifications to improve message delivery, and introduces a RAG-based knowledge Q\&A capability to enhance information access. Test results show that the system is runnable in the current single-node environment: key regression cases passed in reruns, and core request paths such as login, club listing, activity registration, and check-in maintained low response times with a 0\% error rate under the specified concurrency scripts.

Overall, the system can support the major business processes of university club management and demonstrates practical engineering value in functional completeness, authentication, real-time interaction, and auxiliary intelligence. Nevertheless, long-term stability verification, more rigorous AI evaluation, and larger-scale deployment validation still need further improvement.

\vspace{1em}
\noindent\textbf{Keywords:} Spring Boot; Club Management System; JWT Authentication; Real-time Notification; RAG Q\&A

\newpage
